% ============================================================================
% Suggested manuscript additions: EM trajectory and mechanism studies.
% FOR MANUAL REVIEW BY THE AUTHORS.  Nothing here has been applied to
% paper/main.tex, which reproduces the 2026-09-30 manuscript.
%
% Reference manuscript: the authors' latest PDF, received 2026-10-08
% (23 pages; sha256 f38961e0f210b91ba25246fda1ecab2ddfb93988ca7d90d81d73b0520548746b).
% That PDF is not stored in this repository.  Section names and quoted wording
% below come from it.  No \label or \ref to existing manuscript objects is
% assumed; comments mark where a cross-reference belongs.  New labels are left
% to the authors.
%
% Terminology used below.  "EM55" is EM after 55 iterations.  The "endpoint" is
% EM after the iteration cap of 10^6 updates; it is not automatically
% converged, so its certification status is stated where it matters.  "The EM55
% recommendation" is the allocation date chosen from the EM55 fit; that date is
% often LATER than the endpoint's.
%
% Every number is verified from saved outputs by
%   python -m experiments.em_mechanism.verify_claims
% (results/em_mechanism/verification/claims_check.md).
%
% Contents
%   1. Section 7.2 (Decision performance): revised final paragraph
%   2. Section 8 (Conclusion): revised mechanism sentences
%   3. Appendix C.2: replacement for "Robustness to initialization is not evaluated."
%   4. Appendix D.3: replacement for its final sentence
%   5. New appendix subsection (suggested placement: after D.3)
%   6. Figure caption (mechanism figure) and compact evidence table
% ============================================================================


% ----------------------------------------------------------------------------
% 1. SECTION 7.2 "Decision performance", final paragraph.
% PLACEMENT: replaces the existing paragraph that begins
%   "A supplementary trajectory study at tau = 2 suggests that early stopping
%    preserves an upward bias in the estimated third moment ..."
% WHY: the PDF's "77%" and "approximately 89%" hold conditional on the 609
% deteriorating cases (of 660) in which the intervention is feasible; the
% revision states the denominators and the qualifications.
% ----------------------------------------------------------------------------
A supplementary trajectory study at $\tau=2$ points to one channel behind the
gap between early stopping and continued fitting on the reference
populations.  After 55 iterations, EM's estimate of the third moment
$m_3=\mathbb{E}[p^3]$ has a positive median error for every initialization
tested.  Continued fitting lowers it, often below its true value.  On these
populations the welfare loss is asymmetric in this error.  Underestimating
$m_3$ is associated with early intervention and large regret, while moderate
overestimation is associated with much smaller regret.  The upward bias can
therefore help the decision without implying a more accurate estimate of the
risk distribution.  At $N=10^3$, regret rises between iteration 55 and the
iteration cap in 660 of 1{,}500 datasets.  Post-hoc moment interventions
support this explanation and identify the change in $m_3$ as a major
channel.  Keeping the 55-iteration $m_3$ while adopting the later $m_1$ and
$m_2$ preserves the EM55 recommendation in 470 of the 609 cases where the
intervention is feasible (77.2\%; 51 cases are infeasible).  It removes
89.3\% of the additional regret in those cases, or 77.4\% if the infeasible
cases count as unrepaired.  At $N=10^5$ the corresponding figure is 77 of 110
(70.0\%).  The explanation depends on the population: the same bias raises
regret under compatible alternatives whose optimal date is earlier.
% Cross-reference: point to the new appendix subsection (item 5) here.


% ----------------------------------------------------------------------------
% 2. SECTION 8 "Conclusion".
% PLACEMENT: replaces the sentences
%   "The supplementary mechanism study suggests that early stopping preserves a
%    beneficial bias in the estimated third moment, limiting costly shifts
%    toward earlier intervention on the reference populations. This advantage
%    is population-dependent and can reverse under compatible alternatives."
% NOTE: the latest PDF no longer says the mechanism is unresolved; this is a
% tightened version of its current sentences.
% ----------------------------------------------------------------------------
On the reference populations, moment interventions support the explanation
that early stopping retains an upward bias in the estimated third moment.
Under the asymmetric welfare loss observed on this benchmark, that bias
limits costly shifts toward earlier intervention, and continued fitting often
removes it.  The interventions identify this as a major channel, not the only
one.  The advantage depends on the population and can reverse under
compatible alternatives whose optimal date is earlier.


% ----------------------------------------------------------------------------
% 3. APPENDIX C.2 "Estimation and bootstrap aggregation".
% PLACEMENT: replaces the sentence "Robustness to initialization is not
% evaluated." (end of its second paragraph).  The preceding sentence, about the
% reported procedures starting from uniform weights, still holds and stays.
% ----------------------------------------------------------------------------
A separate trajectory study evaluates eight initializations: uniform, three
Beta shapes and four random full-support starts.  The 55-iteration estimate
has lower mean squared regret than certified NPMLE for every one of them on
this benchmark; see the appendix subsection on the trajectory and mechanism
study.
% Cross-reference: point to the new appendix subsection (item 5).


% ----------------------------------------------------------------------------
% 4. APPENDIX D.3, final paragraph ("Results and interpretation").
% PLACEMENT: replaces the closing sentence
%   "These results do not imply improvement on every pair or both members,
%    establish worst-case performance over the full moment class, or explain
%    why early stopping helps."
% ----------------------------------------------------------------------------
These results do not imply improvement on every pair or both members, or
establish worst-case performance over the full moment class.  A separate
trajectory study examines a mechanism behind the early-stopping advantage.
% Cross-reference: point to the new appendix subsection (item 5).


% ----------------------------------------------------------------------------
% 5. NEW APPENDIX SUBSECTION.
% PLACEMENT (suggested): after Appendix D.3, inside Appendix D "Supporting
% decision-performance results", as the next subsection.  It does not exist
% yet; choose its number and label.
% ----------------------------------------------------------------------------
\subsection{Trajectory and mechanism study of early-stopped EM}
% add \label{...} of your choice

\paragraph{Design.}
We record complete EM trajectories at $\tau=2$ on the 200-point grid, with the
same pre-failure likelihood and decision rule as the main benchmark.  Each run
starts from one of eight initializations and stops at an iteration cap of
$10^6$ updates.  At every recorded iteration we store the likelihood, its
directional-derivative optimality residual, the fitted cell probabilities and
moments, the estimated welfare curve and the selected date.  Datasets are
shared across methods: $30$ configurations, $N\in\{10^3,10^5\}$, and $50$
datasets per configuration.  The first 20 datasets per configuration are those
of the main benchmark, and our implementation reproduces the archived 30-, 55-
and 80-iteration decisions on every shared dataset.  We call the state at the
cap the \emph{endpoint}.  The endpoint is not necessarily converged: it meets
the $10^{-8}$ optimality certificate in 1{,}042 of 1{,}500 uniform-start runs
at $N=10^3$ and 1{,}366 of 1{,}500 at $N=10^5$.

\paragraph{Events.}
A dataset \emph{deteriorates} if its regret under the reference population is
higher at the endpoint than after 55 iterations.  This happens in 660 of
1{,}500 datasets at $N=10^3$ and 110 of 1{,}500 at $N=10^5$; regret falls in
32 and 150 datasets, respectively.  In over 92\% of deteriorating datasets the
endpoint recommends an earlier allocation date than EM55.

\paragraph{What changes, and how it differs by sample size.}
From iteration 55 to the endpoint, the fitted $m_3$ falls by a median of about
27\% in deteriorating datasets at both sample sizes; $m_1$ and $m_2$ change
much less.  The two sample sizes otherwise differ.  At $N=10^3$, the fitted
cell probabilities move farther from the true observation probabilities in
94.4\% of deteriorating datasets, and the fitted distribution concentrates on
a few grid points.  At $N=10^5$, they move farther from the true probabilities
in only 20.9\%, and the support barely changes.  Fitting sampling noise
therefore describes the change at $N=10^3$, but not at $N=10^5$.

\paragraph{Prespecified tests.}
Each test solves one linear program for the grid distribution closest in
$L_1$ to a stated target, subject to stated equality constraints, and records
the date it recommends.  The first test fixes the endpoint's fitted cell
probabilities and stays as close as possible to the 55-iteration fit.  That
one constructed distribution recommends the endpoint's date in 98.8\% of
deteriorating datasets at $N=10^3$ and 98.2\% at $N=10^5$.  The test says
nothing about other distributions with the same fitted probabilities.
Freezing unidentified higher moments at their 55-iteration values is almost
never feasible.  At $N=10^3$, freezing the masses of the risk regions that
carry the welfare change preserves the EM55 recommendation in 5.5\% of
feasible cases, no more than freezing irrelevant regions (5.9\%).  These
prespecified analyses do not support an explanation through risk-region
masses or unidentified moments.

\paragraph{Post-hoc moment interventions.}
At $\tau=2$ the fitted cell probabilities are a one-to-one function of
$(m_1,m_2,m_3)$, so we next intervened on these moments separately; the
evidence table below reports the results.  Keeping the 55-iteration $m_3$
while adopting the endpoint's $m_1$ and $m_2$ preserves the EM55
recommendation in 77.2\% of feasible deteriorating datasets at $N=10^3$ (470
of 609; 660 attempted, 51 infeasible) and in 70.0\% at $N=10^5$ (77 of 110).
At $N=10^3$ it removes 89.3\% of the additional regret among feasible cases,
or 77.4\% when infeasible cases count as unrepaired.  Changing $m_1$ or $m_2$
alone rarely moves the recommendation.

\paragraph{Direction of the $m_3$ error.}
After 55 iterations, the median relative error of $m_3$ is between $+5.3\%$
and $+16.8\%$ at $N=10^3$, and between $+3.0\%$ and $+12.6\%$ at $N=10^5$,
across all eight initializations.  This includes one start that begins 80\%
below the truth.  Endpoints and certified NPMLE are approximately
median-unbiased for $m_3$, so their estimate falls below the truth in about
half of the datasets.  Pooled over estimators, regret is much larger after
underestimation than after moderate overestimation.  At $N=10^3$, mean regret
is $0.086$ for errors below $-50\%$ and $0.008$--$0.009$ for errors between
$-10\%$ and $+50\%$; it rises again to $0.022$ for overestimates above
$+50\%$.  Underestimation mostly leads to dates earlier than the optimum.
Replacing the endpoint's $m_3$ with its true value preserves the EM55
recommendation in only 30.6\% of feasible cases at $N=10^3$.  This
explanation concerns the direction of the $m_3$ error under an asymmetric
loss.  It does not compare the overall accuracy of the two estimates of the
risk distribution.

\paragraph{Noise-free comparison.}
With exact observation probabilities (all trajectories certified at a cap of
$10^7$), the recommended date changes between iteration 55 and the endpoint in
33 configuration--initialization cases; 32 of these changes improve welfare
under the reference population and one reduces it.

\paragraph{Limitations.}
The moment interventions were specified after the prespecified tests.
Holding $m_3$ fixed also blocks most improvements (76.7\% and 90.7\% of
improving datasets at $N=10^3$ and $N=10^5$), so it identifies the channel
through which further fitting changes the recommendation, not a protective
rule.  Distribution shape still matters: combining the endpoint's shape with
the 55-iteration moments preserves the EM55 recommendation in only about 44\%
of deteriorating datasets.  Different distributions with identical fitted
probabilities can recommend different dates, but the deterioration studied
here comes with a change in the fitted moments.  The explanation is specific
to these reference populations at $\tau=2$.  The upper member of a compatible
pair has an earlier optimal date in 28 of 30 pairs, and there the same bias
increases regret.


% ----------------------------------------------------------------------------
% 6a. FIGURE (mechanism).
% FILE: results/em_mechanism/figures/fig5_mechanism.pdf (copy into the
% manuscript's figure directory).  PLACEMENT: inside the new subsection.
% ----------------------------------------------------------------------------
\begin{figure}[t]
\centering
\includegraphics[width=\linewidth]{fig5_mechanism.pdf}
\caption{\textbf{One channel behind the early-stopping advantage on the
reference populations ($\tau=2$, uniform grid with 200 points).}
(a) Mean regret against the relative error of the fitted third moment $m_3$,
pooling 55-iteration, endpoint and NPMLE fits.  Underestimation is associated
with much larger regret than moderate overestimation; regret rises again for
overestimates above $+50\%$.
(b) Share of datasets in which $m_3$ is underestimated.
(c) Deteriorating datasets: share of feasible constrained estimates that
preserve the EM55 recommendation.  The first three bars keep the endpoint's
fitted probabilities: closest to the 55-iteration fit, with frozen
decision-relevant risk masses, and with frozen irrelevant masses.  The last
two bars are post-hoc moment interventions (55-iteration $m_3$; true $m_3$).
(d) Mean change in risk mass from iteration 55 to the endpoint at $N=10^3$, for
deteriorating and improving datasets.}
% add \label{...} of your choice
\end{figure}


% ----------------------------------------------------------------------------
% 6b. COMPACT EVIDENCE TABLE.  PLACEMENT: inside the new subsection.
% Requires booktabs (already used by the manuscript's tables).
% ----------------------------------------------------------------------------
\begin{table}[t]
\centering
\small
\caption{\textbf{Interventions on deteriorating datasets.}  Each row solves
one linear program.  Success means the constrained estimate preserves the
EM55 recommendation, except in the row marked $\dagger$, where it means the
estimate moves to the endpoint's recommendation.  Entries are
successful/feasible (share of feasible).  Attempted cases: 660 at $N=10^3$
and 110 at $N=10^5$.  Rows marked $*$ are post hoc.}
% add \label{...} of your choice
\begin{tabular}{@{}lrr@{}}
\toprule
Constrained estimate & $N=10^3$ & $N=10^5$ \\
\midrule
Endpoint fit, closest to iteration 55 & 8/660 (1.2\%) & 2/110 (1.8\%) \\
\quad $+$ frozen decision-relevant risk masses & 7/128 (5.5\%) & 3/75 (4.0\%) \\
\quad $+$ frozen irrelevant risk masses & 8/136 (5.9\%) & 10/96 (10.4\%) \\
\quad $+$ frozen $m_4$--$m_6$ & 1/2 & 0/0 \\
$*$ Endpoint $m_1,m_2$; iteration-55 $m_3$ & 470/609 (77.2\%) & 77/110 (70.0\%) \\
$*$ Endpoint $m_1,m_2$; true $m_3$ & 148/484 (30.6\%) & 24/110 (21.8\%) \\
$*\dagger$ Iteration-55 fit with endpoint $m_3$ & 133/206 (64.6\%) & 28/54 (51.9\%) \\
\midrule
Regret increase removed by iteration-55 $m_3$ & 89.3\% & 84.8\% \\
\quad infeasible cases counted as unrepaired & 77.4\% & 84.8\% \\
\bottomrule
\end{tabular}
\end{table}
